\documentclass[UTF8,a4paper,11pt]{ctexart}
\usepackage[margin=24mm]{geometry}
\usepackage{amsmath,amssymb,booktabs,hyperref}
\hypersetup{colorlinks=true,linkcolor=blue,urlcolor=blue}
\title{受约束轮腿残差PPO的任务一致性与通道效应审计}
\author{作者待补}
\date{2026年10月4日}
\begin{document}
\maketitle
\noindent\textbf{稿件状态：}方法、开发预试验与独立固定模型确认已完成。
本文为可供正式写作和导师审阅的实证审计初稿，作者信息、目标期刊格式与语言润色待补；
不以材料就绪保证录用。内置编译环境不可用，当前源码未确认生成PDF。

\begin{abstract}
输出限幅、回报目标等价与完整任务成功是不同层面的要求。
本文在固定VMC与六状态LQR轮腿系统中审计三维残差PPO，
低层与物理以2 kHz运行，策略50 Hz，保留实际被动链几何和执行器证据。
原奖励与已知势函数奖励各使用三个配对种子、同初始条件和200k连续策略步，
总计1.2M；公开开发结果中塑形成功均值差为$-8.33$个百分点，原继续门未通过。
随后固定全部六个末模型，在事前登记的新常规96、压力48与旧28回归上
完成4472回合的轮/腿通道闭环干预与强经典对照。
新常规中禁用轮残差的六模型成功平均增加62.5个百分点，
但设计失败增加5例；压力与旧能力结果分别报告，提前终止失败不删除。
本研究提供限定系统的任务--执行契约、完整负结果及独立通道效应证据，
不提出新PPO或PBRS算法，不将通道禁用视为同权限训练改进，
也不将输出包络通过解释为全局状态安全。
\end{abstract}
\noindent\textbf{关键词：}轮腿机器人；残差强化学习；约束执行；任务一致性；闭环通道干预

\section{研究问题与主张边界}
经典控制器上叠加学习残差、比较动作表示以及势函数奖励塑形均已有研究。
本文不将PPO、虚拟力矩或势函数的组合视为新算法。
本文检验的问题是：在固定轮腿系统、执行器包络与累积任务判据下，
输出可行性、折扣目标一致性与有限预算的完整任务表现是否相符，
以及学习残差的轮、腿通道在闭环中分别产生何种影响。
已知方法对照包括残差学习、动作空间设计和层次势函数塑形工作\cite{residual,action,hprs}。

本文只讨论仿真地面任务。跳跃、单侧坡演示、ROS2与实机安全不属于本次结果范围。
局部线性设计与扭矩限幅不构成全局稳定性或状态安全不变集证明。

本文的实证贡献限定为三项。
其一，统一报告低层输出、设计域、真实几何、接触任务和评价终态，
使``输出可行''与``完成任务''的差异在同一数据中可核查。
其二，用同初始化、同预算及同信息的三种子学习比较检验奖励收益主张，
保留失败出口与种子异质性，而不以单个有利种子或次指标替换主门。
其三，对固定末控制器登记并执行独立通道干预，
报告常规改善、设计权衡、压力与旧能力边界及可复现源合同。
这些是该系统上的实证证据和审计材料，不是关于所有机器人或所有残差学习的定理。

\section{相关工作与实验区别}
残差控制已用经典反馈与学习补偿叠加处理接触问题\cite{residual}。
Hybrid LMC将LQR与集成SAC用于轮式人形模型\cite{hybrid}，
因此本文不以``LQR加PPO''作为贡献，也不横向比较不同平台的成功率。
动作表示研究已显示初始化及策略步间行为会影响学习\cite{action}；
本文的通道mask不是动作空间的重新训练比较，不能据其优势归因于维数或物理先验。

受约束残差控制已有指定系统假设下的稳定性分析\cite{crrl}。
本文的欠驱动接触切换、构型映射和基控制记忆并未验证满足那些定理前提，
因而只报告实际输出及状态轨迹证据。
HPRS已将安全、目标和舒适需求组织为势函数奖励并采用PPO\cite{hprs}。
本文使用已知势函数作为同折扣目标的受控干预，不重新声称塑形或约束奖励首创。
实验区别在于固定高频残差执行链、延迟物理包与当前请求的混合时序、
累积任务判据，以及固定模型通道作用与强经典控制的共同审计。
这里没有复现所有前述作者模型；单一系统实证不构成排他性创新穷尽证明。

\section{模型、控制与执行契约}
仿真使用MuJoCo Warp。固定名义设计质量为7 kg，训练被控对象参数可变；
名义模型与实际随机对象不得混同。
低层使用VMC与六状态LQR，物理和低层更新周期为0.5 ms，策略周期为20 ms。
每个策略动作保持40个物理子步，但构型映射、执行器余量和限幅逐子步更新。
腿长目标在每回合内固定，覆盖0.115--0.380 m；训练采用每阶段每种子的100个固定场景，
不是每次回合独立重采样参数。

M3策略产生三维差模请求
\begin{equation}
 \bar a=(\bar a_F,\bar a_H,\bar a_W)\in[-1,1]^3,\quad
 \Delta F_L=-\Delta F_R=c_F \bar a_F,\quad
 \Delta H_L=-\Delta H_R=c_H \bar a_H,\quad
 \Delta \tau_{WL}=-\Delta \tau_{WR}=c_W \bar a_W .
\end{equation}
其中$\bar a$表示执行链的滤波请求而非未经过滤的策略样本。
尺度沿用冻结配置：$c_F=0.1\times7\times9.81/2=3.4335$ N，
$c_H=1$ N\,m，$c_W=0.3$ N\,m。
构型相关映射记为$B(q)$。先形成完整名义输出，再处理残差和总输出约束：
\begin{equation}
 \tau=\tau_{\rm nom}+\lambda B(q)\bar a,\qquad 0\leq\lambda\leq1,
 \qquad \tau_i^{\min}(\dot q_i)\leq\tau_i\leq\tau_i^{\max}(\dot q_i).
\end{equation}
名义请求可行化和最终执行器限幅属于实际输出链，不能省略后再讨论残差权限。
由于左右构型、接触和饱和可能不同，虚拟差模并不保证实际运动严格解耦。

真实被动链几何、八个关节余量、实际驱动力矩和名义命令逐物理步记录。
主动关节设计域为1.4 rad，物理限位为1.5 rad，二者是不同判据。
最小链长0.114704466 m是仿真几何代理边界，不是经过实机认证的安全值。

\section{可得观测与完整任务}
Actor和Critic共享38维接口：32维物理包包含姿态、角速度、速度、命令误差、
高度目标、关节、轮速、腿长及实际残差；另6维为当前自身滤波请求。
物理包可延迟0--10 ms，自身请求保持当前时序。
绝对横向位置、完整名义控制器内部记忆与累积约束历史未全部包含。
因此本文不宣称该接口为完整Markov状态，也不将瞬时缺失位置等同于历史无法估计位置。

成功须同时满足完整到达、保持与停车、所需障碍接触和退出、
姿态、速度、目标高度、设计域及物理证据约束。
使用地形相对姿态时相对峰值为5度，世界roll/pitch另有10度上限；yaw峰值仍为5度。
速度RMSE不超过目标速度绝对值的20\%，停距不超过0.6 m，
到达1.5 s后的尾速不超过0.03 m/s，到达后保持2 s。
高度RMSE和终态平均FK腿长误差均不超过0.02 m。
接触位保留原模拟器几何接触语义，不代表正法向载荷认证。
统计失败可重叠，不能把其求和解释为互斥原因。

\section{奖励变换与价值估计边界}
原密集奖励包含速度跟踪、姿态、实际残差幅值及变化率项，逐0.5 ms积分。
完整终态按原成功判据获得$+10$或$-10$。
本研究已完成的奖励预试验使用已知势函数形式\cite{ng,hprs}：
\begin{equation}
 r'_t=r_t+\gamma\Phi_{t+1}-\Phi_t,\qquad \gamma=0.99.
\end{equation}
中间势函数为所定义物理、设计或姿态历史违规时$-10$，否则为0；初态和所有终态设为0。
它没有覆盖全部接触、停车、尾速或高度任务失败。
有限完整回合满足
\begin{equation}
 \sum_{t=0}^{T-1}\gamma^t(r'_t-r_t)
 =-\Phi_0+\gamma^T\Phi_T=0.
\end{equation}
非终态截断不能丢弃右端边界项。检查点保存不重启物理、势函数或RNG；
实际用于上一rollout的bootstrap预测、done和更新前epoch单独记录。
若状态价值同步变换为$V'=V-\Phi$，TD残差亦可保持一致，
但38维近似Critic未必能表达历史势函数。
目标恒等不证明PPO收敛、实际优势或纯信用分配因果。
未折扣的塑形回报不必相同，报告继续使用原任务和原episode回报。

\section{完成的有限学习预试验}
两组分别使用原奖励和势函数奖励，三个配对种子1609、1610、1611。
每组每种子连续训练200k策略步，共1.2M；100世界、50步rollout、batch 250、
10个epoch、2层64单元MLP及标准PPO损失保持一致。
每run实际400个训练epoch、8000次Adam更新；参数、Adam和物理在CUDA运行。
初始权重、世界配置及所有检查点计数核验。
公开开发96案例只使用固定200k末模型，中间检查点不用于选模。

\begin{table}[ht]
\centering
\caption{两奖励组保持一致的学习与观测配置。}
\begin{tabular}{lr}\toprule
配置 & 数值\\\midrule
Actor/Critic隐藏层 & $64,64$\\
学习率 & $3\times10^{-4}$\\
折扣/GAE & $0.99/0.95$\\
PPO clip/价值clip & $0.2$/关闭\\
熵系数/价值系数 & $0/0.5$\\
最大梯度范数 & $0.5$\\
观测归一化/奖励归一化 & 开启/关闭\\
观测clip/epsilon & $10/10^{-8}$\\
M3初始log-std & $(-0.717260,-0.725599,-1.384311)$\\\bottomrule
\end{tabular}
\end{table}
初始log-std来自冻结的物理响应校准，只有新策略创建时设置；
保存恢复不重新初始化。
CUDA用于Actor、Critic、Adam矩及MuJoCo Warp物理，
Python收集与归一化含CPU工作，不将物理设备或吞吐等同端到端加速。

\begin{table}[ht]
\centering
\caption{同一公开开发96案例上的固定末模型结果；不是独立终评。}
\begin{tabular}{lrrrr}\toprule
种子 & 原奖励成功 & 势函数成功 & 原设计失败 & 势函数设计失败\\\midrule
1609 &44&38&3&6\\
1610 &65&6&0&1\\
1611 &7&48&2&1\\\bottomrule
\end{tabular}
\end{table}
势函数成功均值差为$-8.33$个百分点，仅一个种子改善，设计失败多3例。
原预注册继续门未通过，停止该收益分支，不追加2M训练或扫描幅度。
同一案例被不同训练代理重复评价，116/288和92/288不是288个独立随机场景。
三种子的异质性全部保留，不宣称普遍有害或统计显著性。

\section{独立固定模型通道确认协议}
公开探索中，固定势函数seed1610模型在全部残差、仅腿、仅轮、零残差条件下成功数为
6、93、8、93/96。该闭环mask干预改变可达请求子空间及后续状态和动作历史，
仅能解释该固定策略系统的通道效果，不是同权限训练方法比较。

新的确认固定全部六个200k末模型、四个0/1mask及强经典B0/B1。
新常规96、六因素压力48和旧公开28回归分别报告，预算4472回合、训练0步。
主估计为同模型同参数案例的$\operatorname{success}(\text{legs})-
\operatorname{success}(\text{all})$，六模型和三个种子聚类全部报告，不选case或模型。
旧28回归不是独立数据；原64研究门及封存3000不用于该确认或调参。

\textbf{结果状态：}独立固定模型评价已完成，通道效应及失败边界见下节；
相对已有工作的可发表贡献和完整稿件讨论仍待审阅。
经典压力中已有超时失败；物理包络通过不能替代完整任务成功。
未完成到达轨迹的航向分数不补零，也不从统计中删除。

\section{已完成的独立固定模型确认}
独立确认78个job、4472案例终态全部返回并核验，实际训练0步。
新的常规96与压力48在模型及mask固定后生成，旧28回归非新独立数据。
所有输入、源合同、案例参数、物理证据步数、成功重建与清单哈希核验。
22条方法案例提前终止保留在成功率分母中，不补不可计算的航向分数。

\begin{table}[ht]
\centering
\caption{固定模型完整任务成功次数：全部残差与禁用轮残差。}
\begin{tabular}{lrrrrrr}\toprule
&\multicolumn{2}{c}{新常规96}&\multicolumn{2}{c}{压力48}&\multicolumn{2}{c}{旧回归28}\\
模型&全&腿&全&腿&全&腿\\\midrule
Raw 1609 & 30 & 87 & 14 & 35 & 14 & 28\\
Phi 1609 & 36 & 88 & 12 & 35 & 19 & 27\\
Raw 1610 & 55 & 92 & 25 & 35 & 28 & 28\\
Phi 1610 & 4 & 92 & 5 & 35 & 7 & 28\\
Raw 1611 & 5 & 91 & 4 & 35 & 11 & 28\\
Phi 1611 & 50 & 90 & 21 & 35 & 18 & 28\\
\bottomrule\end{tabular}
\end{table}

新常规主估计$\bar d=0.625$；六模型逐一成功增加57、52、37、88、86、40例，
三个种子簇的描述差为56.77、65.10、65.63个百分点。
这是固定有限模型集合的闭环干预效果，不是六次独立训练重复的显著性检验。
常规设计失败从全部残差11例增加至禁轮16例，物理包络通过均为576/576。
压力成功总次数为81/288对210/288，但设计失败也由11增加到14；
旧回归成功为97/168对167/168，两条件物理与设计均为168/168。
这些分母为模型--案例记录，不是288或168个独立随机环境。

强经典B1在新常规成功93/96、压力36/48、旧回归27/28。
禁轮后的学习腿残差均值未显示对强经典的任务优势，零残差结果同样保留。
通道mask不是新学习算法，减少控制请求子空间可能改变腿通道的实际执行余量，
不能将成功恢复直接解释成设计安全提升或学习失败的唯一原因。
本确认没有改变已失败的奖励收益门，也没有使用旧64门或封存3000调参。

\section{统计单位与可复現性}
记模型$m$在固定案例$i$的成功指示为$S_{mi}^{c}$，其中$c$表示通道条件。
确认的主效果按原登记计算：
\begin{equation}
 d_m=\frac{1}{96}\sum_{i=1}^{96}
       \left(S_{mi}^{\rm legs}-S_{mi}^{\rm all}\right),\qquad
 \bar d=\frac{1}{6}\sum_{m=1}^{6}d_m.
\end{equation}
同一训练种子的两种奖励模型为一个种子簇，先报告六个模型效果及三个种子簇，
而不是将六个模型或576次案例评价当作独立训练重复。
常规96是预先生成的有限面板，腿长含固定节点与连续样本，
其重复控制器评价不扩展独立场景数。
压力每因素8案例属于预定义压力点，不作整个连续压力域的覆盖保证。
旧28案例仅检验既有能力，不是新的独立终评。
三种子不足以支持精确的训练分布结论；本稿不以其宣称统计显著性或泛化定理。

表格由核验后的案例数据自动导出，避免手工复制模型或成功次数。
压力汇总行与六个分因素行描述相同48个案例，不能相加得到额外样本。
表格行数也不是评价job或回合数。
只有全部六个常规模型的配对齐全时才生成主估计，
只有4472案例和78个job全部返回且通过复核后才讨论整体确认结果。

成功、物理包络、设计域、姿态、接触和早终止分别列报。
通道干预可能使成功增多而设计失败增多，不能只挑有利指标。
具有完整到达保持轨迹的航向指标与所有案例成功率是不同分母；
未得到完整航向轨迹时保留不可计算标记、终止原因和该案例失败。

复现入口依次为固定学习预试验的提案、训练源合同、初始化记录、
更新后检查点、训练回合记录、固定末模型开发结果以及新的通道确认协议。
确认结果包含每案例参数与种子、原始判据、物理证据步数、
源合同和文件哈希，已完成job另有追加清单。
只读复核保存实际读取的清单字节快照，防止生产者仍在运行时发生清单错配。
模型与归一化工件保留，CPU重读只用于核保存状态；实际学习与物理运行设备单独验证。

\section{讨论与适用范围}
\subsection{输出约束与任务成功的差异}
常规比较中两种主要通道条件的物理证据均通过，
完整成功却分别为180/576和540/576，说明输出包络和真实关节/链长约束
不能代替接触、姿态、停车与目标高度共同定义的任务。
探索性最大退化模型存在更大的车体横向位移和移动阶段航向偏置，
与缺接触证据相符；该关联不证明未观测路线状态就是训练失败的唯一原因。

\subsection{势函数等价不保证有限PPO改善}
完整回合折扣恒等及端点处理经离线和在线核验，
但塑形学习预试验没有满足原继续门。
这不否定已知势函数理论，而是表明固定信息、近似Critic与有限预算的
学习性能不能从回报恒等直接推导。
若同步平移理想状态价值，TD残差可以不变；真实近似价值网络的目标、
可表示性和优化均可能影响结果。
这里没有价值MonteCarlo校准或追加学习因果干预，故不将失败归于其中唯一机制。

\subsection{通道干预的正效果与设计权衡}
独立常规六模型均显示移除轮请求的任务正效果，
压力及旧回归也显示有限案例改善。
轮通道mask减少可达请求子空间，同时改变后续状态及自身请求历史；
它没有重新训练策略，不能视为同权限算法改进。
设计失败在常规和压力分别增加5和3例，
因此不能把任务恢复称为全部安全指标改进。
强经典B1的常规成功仍高于禁轮腿模型均值，
这使本研究成为控制--任务审计而非提出优于经典控制的新策略。

\subsection{限制与外推}
确定性均值回放中的Critic预测不是随机训练策略的MonteCarlo价值目标。
横向车体代理亦不能替代实际轮胎、障碍几何交集和承重证明。
三种子只提供有限训练重复，压力每因素8点并非连续域覆盖。
旧28回归不是新独立测试；仅仿真、固定名义设计与所登记常规/压力域构成外推边界。
独立案例保持完整，但模型来自已有开发流程，
结论只针对固定模型集合，不替代新的随机学习重复或实机验证。
没有将物理吞吐表述成PPO端到端训练加速，也没有用新确认改写旧失败研究门。

\section{结论}
在所登记系统中，回报目标恒等与输出包络合规均不能保证完整接触任务的学习改善。
三种子同预算奖励预试验按原门失败；固定六个末模型的独立通道确认
显示移除轮请求可恢复大量任务，但伴随设计域权衡，仍没有常规胜强经典的证据。
本文据此提供可复现的限定系统审计及失败边界，
建议在报告学习控制效果时共同核对信息时序、实际执行、累积任务和完整失败记录。
后续若研究路线状态、价值历史或其他控制结构，应重新登记目标、信息和独立评价，
不能据本数据宣称新算法收益或全局安全。

\appendix
\section{复现材料入口}
所有入口相对项目根目录，原记录保持不可变。
学习提案和源合同位于
\texttt{wheelleg\_warp/results/paper\_recovery\_20261004/reward\_pilot\_v1/}，
其中\texttt{proposal.json}、\texttt{trainer\_contract.json}和
\texttt{pair\_review.json}对应设计、源冻结与三种子结果。
独立确认材料位于同级\texttt{channel\_validation\_v1/}，包含
\texttt{protocol.json}、\texttt{source\_contract.json}、\texttt{completion.json}、
\texttt{reviewed\_jobs\_snapshot.json}、\texttt{independent\_review.json}及
\texttt{verified\_tables.csv}。
\texttt{plot\_results.py}再生三面板图；学习与轨迹诊断各自保留图脚本及输入SHA。
\texttt{review\_reward\_pilot.py}、\texttt{review\_channel\_validation.py}和
\texttt{report\_channel\_validation.py}为不重训的工件复核入口。

复核原工件不等于新的物理重复。
物理重复须在新输出命名空间登记，保持场景值、版本、初始化、
时间与硬件条件，并重新冻结源码；不得删除原文件来绕过已存在目录保护。
记录含本机路径时应在新协议明确迁移，不静默改写旧协议哈希。
GPU算术不保证跨运行位级轨迹一致，本文也未复现实机。
原CPU基准、旧研究失败门和封存最终参数保留，未被本确认替换。

\begin{thebibliography}{9}
\bibitem{residual} Johannink T, et al. Residual Reinforcement Learning for Robot Control. 2018.
\url{https://arxiv.org/abs/1812.03201}.
\bibitem{action} E\ss{}er J, Margolis G B, Urbann O, Kerner S, Agrawal P.
Action Space Design in Reinforcement Learning for Robot Motor Skills.
PMLR 270:4021--4032, 2025. \url{https://proceedings.mlr.press/v270/esser25a.html}.
\bibitem{ng} Ng A Y, Harada D, Russell S J. Policy Invariance Under Reward Transformations:
Theory and Application to Reward Shaping. ICML, 1999.
\url{https://ai.stanford.edu/~ang/papers/shaping-icml99.pdf}.
\bibitem{hprs} Berducci L, Aguilar E A, Ničković D, Grosu R.
HPRS: Hierarchical Potential-Based Reward Shaping from Task Specifications.
Frontiers in Robotics and AI, 2025. doi:10.3389/frobt.2024.1444188.
\bibitem{hybrid} Baek D, Purushottam A, Ramos J.
Hybrid LMC: Hybrid Learning and Model-based Control for Wheeled Humanoid Robot via Ensemble Deep Reinforcement Learning.
arXiv:2204.03159, 2022. \url{https://arxiv.org/abs/2204.03159}.
\bibitem{crrl} Staessens T, Lefebvre T, Crevecoeur G.
Adaptive control of a mechatronic system using constrained residual reinforcement learning.
arXiv:2110.02566, 2021. \url{https://arxiv.org/abs/2110.02566}.
\end{thebibliography}
\end{document}
